\section{Custom models}
Every model in the Polarization models help page is a particular closed-form solution of the general line-of-sight integral that defines the complex polarization (see the Complex polarization help page). The custom model builder lets you skip the closed form and solve that integral directly, for whatever emissivity and Faraday-depth profile you type in.

Following Sokoloff et al. (1998) / Burn (1966), consider a line of sight going from the far side of the source at $z_{\rm far}$, through its near side at $z_{\rm near}$, to the observer at $z_{\rm obs}$:
\begin{equation}\label{eq:custom_P}
    \displaystyle P(\lambda) = \frac{1}{I_\lambda}\int_{z_{\rm far}}^{z_{\rm near}} j_p(z)\,e^{2i\phi(z)\lambda^2}\,dz, \qquad I_\lambda = \int_{z_{\rm far}}^{z_{\rm near}} j(\lambda,z)\,dz
\end{equation}
\begin{equation}\label{eq:custom_phi}
    \displaystyle \phi(z) = \int_{\max(z,z_{\rm far})}^{z_{\rm obs}} \phi'(z')\,dz'
\end{equation}
$j_p(z)=p\, j(z)\, e^{2i\chi}$ is the polarized emissivity and $j(\lambda,z)$ the total (unpolarized) emissivity at position $z$, the same $j_p$/$j$ from the Complex polarization help page's general volume integral, here reduced to a single line of sight. $\phi(z)$ is the Faraday depth actually accumulated by a photon emitted at $z$ on its way out: since it only travels forward, it is only rotated by material ahead of it ($z'>z$), not behind it, which is why Equation \ref{eq:custom_phi} integrates from $z$ (not $z_{\rm far}$) up to $z_{\rm obs}$.
\\

\subsection{The builder}
Under Models $\to$ Build Custom Model..., you can type mathematical expressions for $j_p(z)$ and $\phi'(z)$ directly, over a fixed normalized line of sight $z\in[-1,1]$ running the way light travels: $z=-1$ at the source's far side, $z=1$ at the observer. Rather than $z_{\rm far}$/$z_{\rm near}$/$z_{\rm obs}$ each being set once for the whole model, $j_p(z)$ and $\phi'(z)$ each carry their own independent [lower, upper] bounds and are forced to zero outside them, so Equation \ref{eq:custom_phi} becomes
\begin{equation}\label{eq:custom_phi_impl}
    \displaystyle \phi(z) = \int_{\max(z,\,p_{\rm lo})}^{p_{\rm hi}} \phi'(z')\,dz'
\end{equation}
where $[p_{\rm lo},p_{\rm hi}]$ and $[j_{\rm lo},j_{\rm hi}]$ (the bounds on $\phi'(z)$ and $j_p(z)$ respectively, both defaulting to $[-1,1]$) need not coincide or even overlap. Where they overlap, different emission depths pick up different amounts of rotation before reaching the observer, i.e., internal, differential Faraday rotation. Where $\phi'(z)$'s range lies entirely beyond $j_p(z)$'s (e.g. $p_{\rm lo}\geq j_{\rm hi}$), every emission point instead sees the exact same total rotation, i.e., a pure external screen. Each bound field usually just holds a plain number, but may instead be an expression referencing a new constant (e.g. $j_{\rm hi}=3w$), which then gets its own slider exactly like a constant used directly in $j_p(z)$/$\phi'(z)$ themselves.

$j_p(z)$ is typed as its own amplitude/EVPA/shape, multiplying an always-present, non-editable $e^{2i\phi(z)\lambda^2}$ phase factor shown to its right. Its default, $p_0\,e^{2i\chi_0}$, reproduces the uniform, unpolarized-shape-independent amplitude and EVPA every built-in model already has: $p_0$ (fractional polarization amplitude) and $\chi_0$ (EVPA) are, like every other constant here, their own always-present sliders, but typed directly into the field rather than applied automatically, so a screen that needs neither to vary with position can simply leave that default as-is. The field may be complex beyond just $\chi_0$ too (e.g. $e^{2i\chi(z)}$), giving the source its own position-dependent intrinsic EVPA. The plain emissivity $j(\lambda,z)$ used for $I_\lambda$'s normalization is not typed separately: it is the same $j_p(z)$ shape with $p_0$ forced to $1$ and the magnitude taken of whatever is left, since $\chi_0$'s phase (or any other position-dependent phase $j_p(z)$ might carry) has magnitude $1$ either way.

$\phi'(z)$ is a Faraday-depth density (proportional to $n_e B_\parallel$, the same role a rotation measure per unit length plays), not the depth itself; $\phi(z)$ is its accumulated integral, Equation \ref{eq:custom_phi_impl}. Its default, $\phi_0$, is the model's third always-present slider (Faraday-depth scale); a typical screen leaves $\phi'(z)$ as $\phi_0$ or scales a bare shape (e.g. $\phi_0 z$) to set the actual rad m$^{-2}$ magnitude.
\\

\subsection{Syntax}
Both fields (and the four bound fields) accept the line-of-sight coordinate $z$, the reserved symbols $\nu$ (frequency, in MHz) and $\lambda$ (wavelength, in m) for a profile that also depends on frequency directly (e.g. $j_p\propto(\nu/1000)^{-0.7}$ folds a spectral index straight into the emissivity, on top of whatever the Spectrum box's own overall $\alpha$ already applies to the total Stokes $I$), the imaginary unit $i$, and any new constant name with an explicit $*$ for products (implicit multiplication is not supported, since a name like $p0$ would otherwise be read as $p\times0$). Leaving either field blank means the constant $1$.

Supported functions are $\sin$, $\cos$, $\tan$, $\arcsin$, $\arccos$, $\arctan$, $\mathrm{atan2}$, $\sinh$, $\cosh$, $\tanh$, $\exp$, $\log$, $\sqrt{\ }$, $|\cdot|$ (abs), $\mathrm{Heaviside}$, $\mathrm{Piecewise}$, $\min$, $\max$, $\mathrm{erf}$, $\mathrm{floor}$, $\mathrm{ceiling}$, $\mathrm{sign}$, plus the constants $\pi$ and $e$. Two more are built specifically for line-of-sight profiles: $\mathrm{delta}(z_0)$, a thin, concentrated feature at $z=z_0$ (e.g. a razor-thin emitting shell or screen), realized internally as a very narrow fixed-width Gaussian rather than a literal point mass, since the latter cannot be evaluated on any numerical grid; and $\mathrm{gaussian}(z_0,\sigma_z)$, the same idea with a width you choose yourself, a plain number, or a new constant name (with its own slider), for either $z_0$ or $\sigma_z$.
\\

\subsection{Discovering constants and previewing}
Clicking "Parse" scans $j_p(z)$, $\phi'(z)$, and the four bound fields for any name that isn't $z$/$\nu$/$\lambda$/$p_0$/$\chi_0$/$\phi_0$, and adds each one as a row of a "Discovered constants" table, alongside the always-present $p_0$/$\chi_0$/$\phi_0$ (and, if the Spectral model dropdown is set to SSA or Thermal, that shape's own $\nu_0$/$T$) shown for reference. Each new row needs a Kind, which auto-fills sensible bounds and a preview value (all three stay editable afterward): Fraction is a $0-100\%$ quantity, like $p_0$; Angle is a $\pm90^\circ$ quantity, like $\chi_0$; Depth is a signed rad m$^{-2}$ quantity on a log scale, like $\phi_0$; Dispersion is an unsigned rad m$^{-2}$ quantity on a log scale, like a Burn-style $\sigma_\phi$; Number is a plain linear quantity with no particular physical units, for a ratio, count, or generic scale factor; and Frequency/Wavelength is a log-scale quantity in MHz/m, for a constant that sets the model's own characteristic spectral scale (e.g. a break frequency inside $j_p(z)$), distinct from the shared turnover the main window's own Spectrum box controls.

Parsing also refreshes a preview of $j_p(z)$ (real and imaginary parts) and $\phi(z)$/$\phi'(z)$ across the line of sight, each constant evaluated at its own Preview value (and, if $j_p(z)$/$\phi'(z)$ reference $\nu$/$\lambda$ directly, at the Preview $\lambda$ box above the plot), letting you check the profile behaves as expected before committing to it. The Spectral model dropdown previews (but does not set) whether the model picks up opacity: choosing SSA or Thermal there shows the genuine per-$z$ absorption factor $e^{-\tau_\lambda(z)}$ that will multiply $j_p(z)$ once this model is actually selected in the main window and the main window's own live Spectrum box is itself set to SSA/Thermal, that live selection is what actually controls it, not this dropdown.
\\

\subsection{Registering and reusing the model}
Clicking "Define Model" closes the window and adds the new model to every model dropdown, ready to visualize and play with exactly like a built-in one; its sliders are $p_0$/$\chi_0$/$\phi_0$ plus every discovered constant, each following the bounds and Kind just chosen. Models $\to$ Edit Custom Model... reopens the currently-selected custom model's own definition for changes, replacing it in place rather than adding a duplicate. A custom model's definition can be written to disk with Export $\to$ Save Model... and brought back in a later session with Import $\to$ Load Model....

Because custom models are solved for numerically, they are not currently available for QU-fitting of real or simulated data, but can still be used to generate simulated measurements (see the Sampling help page) and to explore RM synthesis (see the RM-synthesis help page) on their predicted spectrum.
